% Mathematical modeling section for an IEEE Transactions on Smart Grid manuscript.
% Intended use after the Introduction/Related Work section:
% % Mathematical modeling section for an IEEE Transactions on Smart Grid manuscript.
% Intended use after the Introduction/Related Work section:
% % Mathematical modeling section for an IEEE Transactions on Smart Grid manuscript.
% Intended use after the Introduction/Related Work section:
% % Mathematical modeling section for an IEEE Transactions on Smart Grid manuscript.
% Intended use after the Introduction/Related Work section:
% \input{sections/02_mathematical_model}

\section{Mathematical Model and Problem Formulation}
\label{sec:mathematical_model}

This section formalizes the distribution restoration problem, the GRL controller observation model, the observation-level false-data injection attack (FDIA) model, and the detection/mitigation metrics used in the IEEE 34-bus and IEEE 123-bus studies. The formulation is intentionally aligned with the experimental workflow: the attacker corrupts only the cyber observation received by the controller, whereas the clean, attacked, and mitigated actions are evaluated on the same physical OpenDSS feeder state.

\subsection{Distribution Feeder and Physical State}
\label{subsec:feeder_model}

Let the distribution feeder be represented by a graph
\begin{equation}
    \mathcal{G} = (\mathcal{N},\mathcal{E}),
\end{equation}
where $\mathcal{N}$ is the set of buses and $\mathcal{E}$ is the set of distribution lines, switches, transformers, and other branch elements represented in the controller graph. For an unbalanced distribution feeder, each bus $i \in \mathcal{N}$ may have an active phase set $\Phi_i \subseteq \{a,b,c\}$, and each edge $e=(i,j) \in \mathcal{E}$ may have a phase set $\Phi_e$.

The physical feeder state at restoration step $t$ is denoted by
\begin{equation}
    x_t =
    \left(
        V_t,\,
        I_t,\,
        P_t,\,
        Q_t,\,
        z_t,\,
        \ell_t,\,
        \omega
    \right),
\end{equation}
where $V_t$ is the vector of bus voltages, $I_t$ is the vector of branch-current or branch-flow quantities, $P_t$ and $Q_t$ are real and reactive power quantities, $z_t$ contains switch and topology status variables, $\ell_t$ contains load service or controllable-load status variables, and $\omega$ denotes the outage scenario, including outaged edges, loadshape information, and other exogenous conditions. In the experiments, $x_t$ is not directly optimized in closed form; it is advanced and evaluated through OpenDSS.

The restoration action at step $t$ is a binary vector
\begin{equation}
    a_t \in \{0,1\}^{n_a},
\end{equation}
where entries correspond to controllable switch and load-restoration decisions exposed by the restoration environment. Let $\mathcal{A}(x_t)$ denote the physically feasible action set under the true feeder state. Feasibility includes operational constraints such as topology consistency, unavailable components, voltage and thermal feasibility, and any environment-specific action constraints.

The physical transition is written abstractly as
\begin{equation}
    x_{t+1} = f_{\mathrm{DSS}}(x_t,a_t,\omega_t),
    \label{eq:physical_transition}
\end{equation}
where $f_{\mathrm{DSS}}(\cdot)$ denotes the OpenDSS-backed transition and power-flow evaluation. Equation \eqref{eq:physical_transition} emphasizes that the physical feeder state evolves according to the true outage scenario and selected action, not according to the attacker's corrupted observation.

\subsection{Restoration Objective and Energy Metric}
\label{subsec:restoration_objective}

Let $p^{\mathrm{srv}}_{i,t}(x_t)$ denote the real power served at load bus or load component $i$ during restoration step $t$. For a restoration horizon of $T$ steps and step duration $\Delta t$, the restored energy of a trajectory $\tau = (x_0,a_0,x_1,\ldots,x_T)$ is
\begin{equation}
    E(\tau)
    =
    \Delta t
    \sum_{t=0}^{T-1}
    \sum_{i \in \mathcal{L}}
    p^{\mathrm{srv}}_{i,t}(x_t),
    \label{eq:restored_energy}
\end{equation}
where $\mathcal{L}$ is the set of load components included in the restoration metric. The experiments report $E(\tau)$ in kWh.

A restoration policy seeks to maximize expected restored energy, possibly with penalties for voltage violations, thermal violations, invalid actions, or switching costs. A generic policy objective can be written as
\begin{equation}
    \max_{\pi}
    J(\pi)
    =
    \mathbb{E}_{\omega}
    \left[
        \sum_{t=0}^{T-1}
        r(x_t,a_t,\omega_t)
    \right],
    \qquad
    a_t \sim \pi(\cdot \mid o_t),
    \label{eq:policy_objective}
\end{equation}
where $r(\cdot)$ is the controller training reward and $o_t$ is the controller observation. The paper's evaluation metric is not the training reward itself; it is restored-energy degradation computed from \eqref{eq:restored_energy}.

\subsection{Controller Observation and GRL Policy}
\label{subsec:observation_policy}

The restoration controller does not observe the full physical state $x_t$ directly. Instead, it receives a structured observation
\begin{equation}
    o_t = h(x_t)
    =
    \left(
        X_t,\,
        B_t,\,
        g_t,\,
        m_t
    \right),
    \label{eq:observation}
\end{equation}
where $X_t$ contains node features, $B_t$ contains edge or branch features, $g_t$ contains global restoration scalars, and $m_t \in \{0,1\}^{n_a}$ is an action/status availability mask. In the studied controller, the direct observation families include bus-voltage measurements, branch-flow measurements, topology/status features, action/status masks, and global scalar features.

The action mask defines a controller-level availability set
\begin{equation}
    \mathcal{A}_m(o_t)
    =
    \left\{
        a \in \{0,1\}^{n_a}
        \;:\;
        a_j \leq m_{t,j},
        \; j=1,\ldots,n_a
    \right\},
    \label{eq:mask_feasible_set}
\end{equation}
for action components whose convention is that $m_{t,j}=0$ marks action $j$ unavailable. The exact environment action semantics may include additional constraints, but \eqref{eq:mask_feasible_set} captures why action/status FDIA can be highly consequential: a false-unavailable mask entry can remove a useful restoration action before the learned policy selects an action.

A trained GRL policy is denoted by
\begin{equation}
    a_t^{0} = \pi_{\theta}(o_t),
    \label{eq:clean_action}
\end{equation}
where $a_t^{0}$ is the clean action selected from the uncorrupted observation. In deterministic evaluation, $\pi_{\theta}$ denotes the policy's deployed action-selection rule. For companion controller studies, the same notation also applies to non-GRL policies such as MLP controllers, with a different parameterization of $\pi_{\theta}$.

\subsection{Observation-Level FDIA Model}
\label{subsec:fdia_model}

An observation-level FDIA is modeled as an attack operator
\begin{equation}
    \tilde{o}_t = \mathcal{A}_{\eta}(o_t),
    \label{eq:attack_operator}
\end{equation}
where $\eta$ specifies the attack family, target observation indices, and attack magnitude or replacement rule. The attacked action is
\begin{equation}
    a_t^{\eta} = \pi_{\theta}(\tilde{o}_t).
    \label{eq:attacked_action}
\end{equation}
The physical feeder state remains $x_t$. Thus, $\tilde{o}_t$ affects the controller's selected action, but the action is evaluated on the true feeder through \eqref{eq:physical_transition}.

Let $\mathcal{S}_{\eta}$ be the set of observation channels modified by the attack. The attacks studied in the evidence packages are structured as follows.

\subsubsection{Action/Status Mask Attack}

A false-unavailable action-mask attack changes one or more availability-mask entries from available to unavailable:
\begin{equation}
    \tilde{m}_{t,j}
    =
    \begin{cases}
        0, & j \in \mathcal{S}_{\eta}^{m},\\
        m_{t,j}, & j \notin \mathcal{S}_{\eta}^{m}.
    \end{cases}
    \label{eq:action_mask_attack}
\end{equation}
This attack is a direct controller-observation attack because the action mask is consumed by the restoration policy. It is especially important when target $j$ corresponds to a load-restoration or switching action that is critical for reconnecting high-value load.

\subsubsection{Topology/Status Attack}

Let $A_t$ denote a topology or adjacency-status representation used in the graph observation. A false-open topology attack sets selected observed edges to open:
\begin{equation}
    \tilde{A}_{t,ij} = 0,
    \qquad
    (i,j) \in \mathcal{S}_{\eta}^{A},
    \label{eq:false_open_attack}
\end{equation}
whereas a false-closed hidden-outage attack sets selected observed edges to closed:
\begin{equation}
    \tilde{A}_{t,ij} = 1,
    \qquad
    (i,j) \in \mathcal{S}_{\eta}^{A}.
    \label{eq:false_closed_attack}
\end{equation}
The physical edge status in OpenDSS is not changed by \eqref{eq:false_open_attack} or \eqref{eq:false_closed_attack}; only the graph observation received by the controller is changed.

\subsubsection{Voltage Measurement Attack}

Let $v_{i,\phi,t}$ be the per-unit voltage magnitude for bus $i$ and phase $\phi$. A false-low voltage attack replaces selected active-phase voltage entries with a reported value $v^{\mathrm{low}}$:
\begin{equation}
    \tilde{v}_{i,\phi,t}
    =
    v^{\mathrm{low}},
    \qquad
    (i,\phi) \in \mathcal{S}_{\eta}^{V}.
    \label{eq:voltage_attack}
\end{equation}
The experiments use this structure to test whether the policy changes restoration decisions when the cyber observation reports artificially low voltage.

\subsubsection{Branch-Flow Hiding Attack}

Let $b_{e,t}$ be the branch-flow or branch-loading feature for edge $e$. A branch-flow hiding attack suppresses selected branch-flow observations:
\begin{equation}
    \tilde{b}_{e,t}
    =
    \alpha_b b_{e,t},
    \qquad
    e \in \mathcal{S}_{\eta}^{B},
    \label{eq:branchflow_attack}
\end{equation}
where $\alpha_b = 0$ gives full hiding. IEEE 123 results show that coordinated branch-flow hiding can be damaging even when individual branch-flow targets are weak.

\subsubsection{Global Scalar Attack}

For a global scalar feature $g_{r,t}$, an additive or multiplicative scalar attack can be written as
\begin{equation}
    \tilde{g}_{r,t}
    =
    \alpha_g g_{r,t} + \beta_g,
    \qquad
    r \in \mathcal{S}_{\eta}^{g}.
    \label{eq:global_scalar_attack}
\end{equation}
In the current evidence, these attacks are included only for scalar observations actually available to the controller.

\subsection{Clean, Attacked, and Mitigated Trajectories}
\label{subsec:trajectories}

For an outage scenario $\omega$, define the clean trajectory
\begin{equation}
    \tau^{0}(\omega)
    =
    \left(
        x^{0}_0,a^{0}_0,x^{0}_1,\ldots,x^{0}_T
    \right),
    \qquad
    a^{0}_t = \pi_{\theta}(h(x^{0}_t)).
    \label{eq:clean_trajectory}
\end{equation}
The attacked trajectory is
\begin{equation}
    \tau^{\eta}(\omega)
    =
    \left(
        x^{\eta}_0,a^{\eta}_0,x^{\eta}_1,\ldots,x^{\eta}_T
    \right),
    \qquad
    a^{\eta}_t = \pi_{\theta}(\mathcal{A}_{\eta}(h(x^{\eta}_t))).
    \label{eq:attacked_trajectory}
\end{equation}
The mitigated trajectory is
\begin{equation}
    \tau^{M}(\omega)
    =
    \left(
        x^{M}_0,a^{M}_0,x^{M}_1,\ldots,x^{M}_T
    \right),
    \qquad
    a^{M}_t = \mu_{\psi}(h(x^{M}_t),\tilde{o}^{M}_t,d_t),
    \label{eq:mitigated_trajectory}
\end{equation}
where $\mu_{\psi}$ is the defense action-selection rule, $\tilde{o}^{M}_t$ is the potentially attacked or sanitized observation available to the mitigation layer, and $d_t \in \{0,1\}$ is the detector decision. In the experiments, the clean, attacked, and mitigated trajectories for the same case are evaluated under the same outage-edge scenario and loadshape schedule.

\subsection{Physical Degradation and Recovery Metrics}
\label{subsec:metrics}

Let
\begin{equation}
    E^0 = E(\tau^0), \qquad
    E^{\eta} = E(\tau^{\eta}), \qquad
    E^{M} = E(\tau^{M})
\end{equation}
be the clean, attacked, and mitigated restored energies for the same scenario. The restored-energy loss caused by the attack is
\begin{equation}
    L^{\eta} = E^0 - E^{\eta}.
    \label{eq:loss_raw}
\end{equation}
Because an attack can occasionally improve the restoration outcome by chance, the paper reports positive physical degradation as
\begin{equation}
    D^{\eta}
    =
    \max\left(0, E^0 - E^{\eta}\right).
    \label{eq:physical_degradation}
\end{equation}
The corresponding accidental improvement is
\begin{equation}
    I^{\eta}
    =
    \max\left(0, E^{\eta} - E^0\right).
    \label{eq:accidental_improvement}
\end{equation}

Residual degradation after mitigation is
\begin{equation}
    R^{M}
    =
    \max\left(0, E^0 - E^{M}\right).
    \label{eq:residual_degradation}
\end{equation}
The degradation recovered by mitigation is
\begin{equation}
    G^{M}
    =
    D^{\eta} - R^{M}.
    \label{eq:degradation_recovered}
\end{equation}
For $D^{\eta} > \epsilon$, the clipped recovery ratio is
\begin{equation}
    \rho^{M}
    =
    \operatorname{clip}
    \left(
        \frac{D^{\eta}-R^{M}}{D^{\eta}},
        0,
        1
    \right),
    \label{eq:recovery_ratio}
\end{equation}
and $\rho^{M}=0$ when $D^{\eta}\leq \epsilon$. The clipping prevents accidental over-recovery or numerical artifacts from being interpreted as more than full recovery.

For a set of cases $\Omega$ and attack family $\eta$, the mean degradation, total degradation, and damaging-case rate are
\begin{align}
    \bar{D}_{\eta}
    &=
    \frac{1}{|\Omega|}
    \sum_{\omega \in \Omega}
    D^{\eta}(\omega),\\
    D_{\eta}^{\mathrm{tot}}
    &=
    \sum_{\omega \in \Omega}
    D^{\eta}(\omega),\\
    r_{\eta}^{\mathrm{dmg}}
    &=
    \frac{1}{|\Omega|}
    \sum_{\omega \in \Omega}
    \mathbf{1}
    \left[
        D^{\eta}(\omega) > \epsilon
    \right].
    \label{eq:aggregate_metrics}
\end{align}

\subsection{Detection Model}
\label{subsec:detection_model}

The detector maps an observation and optional reference information to a binary decision:
\begin{equation}
    d_t = \mathcal{D}_{\gamma}(\tilde{o}_t,\mathcal{R}_t)
    \in \{0,1\},
    \label{eq:detector}
\end{equation}
where $\gamma$ denotes the detector profile and $\mathcal{R}_t$ denotes available reference or consistency information. A generic residual detector computes channel residuals
\begin{equation}
    r_{k,t}
    =
    \left|
        \tilde{y}_{k,t}
        -
        \hat{y}_{k,t}
    \right|,
    \qquad
    k \in \mathcal{K},
    \label{eq:residuals}
\end{equation}
where $\tilde{y}_{k,t}$ is an observed channel and $\hat{y}_{k,t}$ is a reference, predicted, or consistency-implied value. The normalized detector score can be written as
\begin{equation}
    s_t
    =
    \max_{k \in \mathcal{K}}
    \frac{r_{k,t}}{\tau_k},
    \label{eq:detector_score}
\end{equation}
with thresholds $\tau_k>0$. The residual detector flags when
\begin{equation}
    d_t^{\mathrm{res}}
    =
    \mathbf{1}
    \left[
        s_t > 1
    \right].
    \label{eq:residual_detector}
\end{equation}

Status-integrity checks are used for binary action/status and topology/status observations. Let $\mathcal{B}$ be a set of binary channels and $\mathcal{C}_{\mathrm{stat}}$ be the set of expected logical consistency constraints among status, topology, and action-mask entries. A generic status-integrity detector is
\begin{equation}
    d_t^{\mathrm{stat}}
    =
    \mathbf{1}
    \left[
        \exists c \in \mathcal{C}_{\mathrm{stat}}
        :
        c(\tilde{o}_t) = 0
    \right].
    \label{eq:status_detector}
\end{equation}
The strict profile uses stronger status-integrity assumptions, while the relaxed profile weakens or removes selected binary-status checks to stress-test stealth behavior.

The deployed detector decision is represented as
\begin{equation}
    d_t
    =
    d_t^{\mathrm{res}}
    \lor
    d_t^{\mathrm{stat}}
    \lor
    d_t^{\mathrm{learn}},
    \label{eq:fused_detector}
\end{equation}
where $d_t^{\mathrm{learn}}$ is included only in learning-detector ablations. The learning detector does not replace the physical residual/status checks in the main evidence; it is evaluated as an additional detection signal.

For a case-level attack, detection is reported as
\begin{equation}
    d(\omega)
    =
    \mathbf{1}
    \left[
        \max_{t=0,\ldots,T-1} d_t = 1
    \right].
    \label{eq:case_detection}
\end{equation}
The detection rate for attack family $\eta$ is
\begin{equation}
    r_{\eta}^{\mathrm{det}}
    =
    \frac{1}{|\Omega|}
    \sum_{\omega \in \Omega}
    d(\omega).
    \label{eq:detection_rate}
\end{equation}
Detection on damaging cases is computed by restricting the denominator to cases with $D^{\eta}(\omega)>\epsilon$.

\subsection{Bypass and Stealth Metrics}
\label{subsec:bypass_model}

A high-impact bypass occurs when an attack is physically damaging but is not detected:
\begin{equation}
    b^{\eta}(\omega)
    =
    \mathbf{1}
    \left[
        D^{\eta}(\omega) > \epsilon_D
        \;\land\;
        d(\omega)=0
    \right],
    \label{eq:bypass_indicator}
\end{equation}
where $\epsilon_D$ is the degradation threshold used to define high-impact damage. The bypass rate is
\begin{equation}
    r_{\eta}^{\mathrm{bypass}}
    =
    \frac{1}{|\Omega|}
    \sum_{\omega \in \Omega}
    b^{\eta}(\omega).
    \label{eq:bypass_rate}
\end{equation}
For stealth-stress experiments, the same structure is used with attack-specific stealth intent. A stealth success is a case in which the attack remains below the chosen detection criterion while retaining the intended physical effect.

\subsection{Mitigation as Candidate Action Screening}
\label{subsec:mitigation_model}

After detection, mitigation is formulated as candidate action screening. Let
\begin{equation}
    \mathcal{C}_t
    =
    \left\{
        a_t^{(1)},a_t^{(2)},\ldots,a_t^{(q)}
    \right\}
    \label{eq:candidate_set}
\end{equation}
be a set of candidate restoration actions. Candidate sources may include:
\begin{align}
    a_t^{\eta} &= \pi_{\theta}(\tilde{o}_t),
    &&\text{attacked-observation policy action},\\
    a_t^{\mathrm{mit}} &= \pi_{\theta}(o_t^{\mathrm{mit}}),
    &&\text{detector-mitigated observation action},\\
    a_t^{\mathrm{ref}} &= \pi_{\theta}(o_t^{\mathrm{ref}}),
    &&\text{trusted-reference observation action}.
    \label{eq:candidate_sources}
\end{align}
The trusted-reference candidate is a robust benchmark candidate in the current evidence. It should not be described as a free replacement for detection, because it assumes access to a reference or screened observation source.

Given a candidate action $a \in \mathcal{C}_t$, define a screening score
\begin{equation}
    \widehat{Q}_t(a)
    =
    \widehat{E}_t(a)
    -
    \lambda_V \widehat{\nu}_t(a)
    -
    \lambda_S \widehat{s}_t(a),
    \label{eq:screening_score}
\end{equation}
where $\widehat{E}_t(a)$ is the estimated restored-energy contribution or rollout value, $\widehat{\nu}_t(a)$ is an estimated voltage or constraint-violation penalty, $\widehat{s}_t(a)$ is an optional switching or operational penalty, and $\lambda_V,\lambda_S \geq 0$ are weights. In the experiments, candidate quality is evaluated using the same physical scenario for clean, attacked, and mitigated comparisons; in an online deployment, \eqref{eq:screening_score} would be computed from a trusted state estimator, digital twin, or reference-screened model.

The selected mitigated action is
\begin{equation}
    a_t^{M}
    =
    \arg\max_{a \in \mathcal{C}_t}
    \widehat{Q}_t(a).
    \label{eq:mitigated_action}
\end{equation}

\subsubsection{Gated Candidate Screening}

In gated screening, mitigation is applied only if the detector gate is active:
\begin{equation}
    a_t^{M}
    =
    \begin{cases}
        \displaystyle
        \arg\max_{a \in \mathcal{C}_t^{\mathrm{gate}}}
        \widehat{Q}_t(a),
        & d_t=1,\\[2mm]
        a_t^{\eta},
        & d_t=0.
    \end{cases}
    \label{eq:gated_screening}
\end{equation}
This mode can fail when the detector misses a damaging attack or when the gated candidate set does not include a sufficiently good reference action.

\subsubsection{Always-Reference Screening}

In always-reference screening, the candidate set includes a trusted-reference action regardless of whether the detector gate is active:
\begin{equation}
    \mathcal{C}_t^{\mathrm{ref}}
    =
    \left\{
        a_t^{\eta},
        a_t^{\mathrm{mit}},
        a_t^{\mathrm{ref}}
    \right\},
    \qquad
    a_t^{M}
    =
    \arg\max_{a \in \mathcal{C}_t^{\mathrm{ref}}}
    \widehat{Q}_t(a).
    \label{eq:always_reference_screening}
\end{equation}
The IEEE 34 evidence uses this as the strongest robust-mitigation benchmark for high-impact action/status attacks.

\subsubsection{Horizon-Screened Mitigation}

For IEEE 123 advanced mitigation, candidate screening can be extended over a horizon. Let $\hat{x}_{t+h}^{a}$ be the predicted rollout state when candidate action $a$ is selected and subsequent actions follow a specified rollout rule. The horizon-screening value is
\begin{equation}
    \widehat{Q}_{t:t+H}(a)
    =
    \sum_{h=0}^{H-1}
    \left(
        \Delta t
        \sum_{i \in \mathcal{L}}
        \hat{p}^{\mathrm{srv}}_{i,t+h}(\hat{x}_{t+h}^{a})
        -
        \lambda_V
        \widehat{\nu}_{t+h}(\hat{x}_{t+h}^{a})
    \right),
    \label{eq:horizon_screening_score}
\end{equation}
and the horizon-screened action is
\begin{equation}
    a_t^{M,H}
    =
    \arg\max_{a \in \mathcal{C}_t}
    \widehat{Q}_{t:t+H}(a).
    \label{eq:horizon_screening_action}
\end{equation}
The variants `horizon\_screen\_raw\_mitigated' and `horizon\_screen\_raw\_mitigated\_reference' differ in whether the candidate pool includes only raw and mitigated actions or also trusted-reference actions.

\subsection{Observation Criticality and Coordinated Attacks}
\label{subsec:criticality_model}

For an observation target $j$ and attack family $\eta$, the single-target criticality score is the mean physical degradation over deterministic cases:
\begin{equation}
    C_{\eta}(j)
    =
    \frac{1}{|\Omega_j|}
    \sum_{\omega \in \Omega_j}
    D^{\eta}_{j}(\omega),
    \label{eq:criticality_score}
\end{equation}
where $\Omega_j$ is the set of cases in which the attack is applicable to target $j$. Targets are ranked by $C_{\eta}(j)$, maximum degradation, damaging-case count, and mitigation residuals.

For a coordinated attack on target set $\mathcal{J}$,
\begin{equation}
    \tilde{o}_t
    =
    \mathcal{A}_{\eta,\mathcal{J}}(o_t),
    \qquad
    \mathcal{J} = \{j_1,j_2,\ldots,j_k\}.
    \label{eq:coordinated_attack}
\end{equation}
The coordinated degradation is
\begin{equation}
    D^{\eta}_{\mathcal{J}}(\omega)
    =
    \max
    \left(
        0,
        E^0(\omega)
        -
        E^{\eta,\mathcal{J}}(\omega)
    \right).
    \label{eq:coordinated_degradation}
\end{equation}

For the IEEE 123 branch-flow subset ablation, let $\mathcal{J}_k$ be the top-$k$ selected branch-flow target subset and let $D_{\mathrm{base}}^{\mathrm{bf}}(\omega)$ be the baseline branch-flow hiding degradation. The total baseline reproduction ratio is
\begin{equation}
    \rho_k^{\mathrm{bf}}
    =
    \frac{
        \sum_{\omega \in \Omega_{\mathrm{bf}}}
        D^{\mathrm{bf}}_{\mathcal{J}_k}(\omega)
    }{
        \sum_{\omega \in \Omega_{\mathrm{bf}}}
        D_{\mathrm{base}}^{\mathrm{bf}}(\omega)
    },
    \label{eq:branchflow_reproduction}
\end{equation}
when the denominator is positive. This metric quantifies whether baseline branch-flow vulnerability is driven by one observation or by a coordinated subset of observations.

\subsection{Scope of the Mathematical Claims}
\label{subsec:model_scope}

The model supports claims for direct attacks on observations available to the current controller. It does not, by itself, support direct load-demand FDIA, DER-availability FDIA, or explicit fault-status FDIA unless those variables are added to the controller observation $o_t$ and evaluated with the same clean/attacked/mitigated physical workflow. This distinction is important for preserving the audit boundary of the IEEE 34-bus and IEEE 123-bus evidence packages.


\section{Mathematical Model and Problem Formulation}
\label{sec:mathematical_model}

This section formalizes the distribution restoration problem, the GRL controller observation model, the observation-level false-data injection attack (FDIA) model, and the detection/mitigation metrics used in the IEEE 34-bus and IEEE 123-bus studies. The formulation is intentionally aligned with the experimental workflow: the attacker corrupts only the cyber observation received by the controller, whereas the clean, attacked, and mitigated actions are evaluated on the same physical OpenDSS feeder state.

\subsection{Distribution Feeder and Physical State}
\label{subsec:feeder_model}

Let the distribution feeder be represented by a graph
\begin{equation}
    \mathcal{G} = (\mathcal{N},\mathcal{E}),
\end{equation}
where $\mathcal{N}$ is the set of buses and $\mathcal{E}$ is the set of distribution lines, switches, transformers, and other branch elements represented in the controller graph. For an unbalanced distribution feeder, each bus $i \in \mathcal{N}$ may have an active phase set $\Phi_i \subseteq \{a,b,c\}$, and each edge $e=(i,j) \in \mathcal{E}$ may have a phase set $\Phi_e$.

The physical feeder state at restoration step $t$ is denoted by
\begin{equation}
    x_t =
    \left(
        V_t,\,
        I_t,\,
        P_t,\,
        Q_t,\,
        z_t,\,
        \ell_t,\,
        \omega
    \right),
\end{equation}
where $V_t$ is the vector of bus voltages, $I_t$ is the vector of branch-current or branch-flow quantities, $P_t$ and $Q_t$ are real and reactive power quantities, $z_t$ contains switch and topology status variables, $\ell_t$ contains load service or controllable-load status variables, and $\omega$ denotes the outage scenario, including outaged edges, loadshape information, and other exogenous conditions. In the experiments, $x_t$ is not directly optimized in closed form; it is advanced and evaluated through OpenDSS.

The restoration action at step $t$ is a binary vector
\begin{equation}
    a_t \in \{0,1\}^{n_a},
\end{equation}
where entries correspond to controllable switch and load-restoration decisions exposed by the restoration environment. Let $\mathcal{A}(x_t)$ denote the physically feasible action set under the true feeder state. Feasibility includes operational constraints such as topology consistency, unavailable components, voltage and thermal feasibility, and any environment-specific action constraints.

The physical transition is written abstractly as
\begin{equation}
    x_{t+1} = f_{\mathrm{DSS}}(x_t,a_t,\omega_t),
    \label{eq:physical_transition}
\end{equation}
where $f_{\mathrm{DSS}}(\cdot)$ denotes the OpenDSS-backed transition and power-flow evaluation. Equation \eqref{eq:physical_transition} emphasizes that the physical feeder state evolves according to the true outage scenario and selected action, not according to the attacker's corrupted observation.

\subsection{Restoration Objective and Energy Metric}
\label{subsec:restoration_objective}

Let $p^{\mathrm{srv}}_{i,t}(x_t)$ denote the real power served at load bus or load component $i$ during restoration step $t$. For a restoration horizon of $T$ steps and step duration $\Delta t$, the restored energy of a trajectory $\tau = (x_0,a_0,x_1,\ldots,x_T)$ is
\begin{equation}
    E(\tau)
    =
    \Delta t
    \sum_{t=0}^{T-1}
    \sum_{i \in \mathcal{L}}
    p^{\mathrm{srv}}_{i,t}(x_t),
    \label{eq:restored_energy}
\end{equation}
where $\mathcal{L}$ is the set of load components included in the restoration metric. The experiments report $E(\tau)$ in kWh.

A restoration policy seeks to maximize expected restored energy, possibly with penalties for voltage violations, thermal violations, invalid actions, or switching costs. A generic policy objective can be written as
\begin{equation}
    \max_{\pi}
    J(\pi)
    =
    \mathbb{E}_{\omega}
    \left[
        \sum_{t=0}^{T-1}
        r(x_t,a_t,\omega_t)
    \right],
    \qquad
    a_t \sim \pi(\cdot \mid o_t),
    \label{eq:policy_objective}
\end{equation}
where $r(\cdot)$ is the controller training reward and $o_t$ is the controller observation. The paper's evaluation metric is not the training reward itself; it is restored-energy degradation computed from \eqref{eq:restored_energy}.

\subsection{Controller Observation and GRL Policy}
\label{subsec:observation_policy}

The restoration controller does not observe the full physical state $x_t$ directly. Instead, it receives a structured observation
\begin{equation}
    o_t = h(x_t)
    =
    \left(
        X_t,\,
        B_t,\,
        g_t,\,
        m_t
    \right),
    \label{eq:observation}
\end{equation}
where $X_t$ contains node features, $B_t$ contains edge or branch features, $g_t$ contains global restoration scalars, and $m_t \in \{0,1\}^{n_a}$ is an action/status availability mask. In the studied controller, the direct observation families include bus-voltage measurements, branch-flow measurements, topology/status features, action/status masks, and global scalar features.

The action mask defines a controller-level availability set
\begin{equation}
    \mathcal{A}_m(o_t)
    =
    \left\{
        a \in \{0,1\}^{n_a}
        \;:\;
        a_j \leq m_{t,j},
        \; j=1,\ldots,n_a
    \right\},
    \label{eq:mask_feasible_set}
\end{equation}
for action components whose convention is that $m_{t,j}=0$ marks action $j$ unavailable. The exact environment action semantics may include additional constraints, but \eqref{eq:mask_feasible_set} captures why action/status FDIA can be highly consequential: a false-unavailable mask entry can remove a useful restoration action before the learned policy selects an action.

A trained GRL policy is denoted by
\begin{equation}
    a_t^{0} = \pi_{\theta}(o_t),
    \label{eq:clean_action}
\end{equation}
where $a_t^{0}$ is the clean action selected from the uncorrupted observation. In deterministic evaluation, $\pi_{\theta}$ denotes the policy's deployed action-selection rule. For companion controller studies, the same notation also applies to non-GRL policies such as MLP controllers, with a different parameterization of $\pi_{\theta}$.

\subsection{Observation-Level FDIA Model}
\label{subsec:fdia_model}

An observation-level FDIA is modeled as an attack operator
\begin{equation}
    \tilde{o}_t = \mathcal{A}_{\eta}(o_t),
    \label{eq:attack_operator}
\end{equation}
where $\eta$ specifies the attack family, target observation indices, and attack magnitude or replacement rule. The attacked action is
\begin{equation}
    a_t^{\eta} = \pi_{\theta}(\tilde{o}_t).
    \label{eq:attacked_action}
\end{equation}
The physical feeder state remains $x_t$. Thus, $\tilde{o}_t$ affects the controller's selected action, but the action is evaluated on the true feeder through \eqref{eq:physical_transition}.

Let $\mathcal{S}_{\eta}$ be the set of observation channels modified by the attack. The attacks studied in the evidence packages are structured as follows.

\subsubsection{Action/Status Mask Attack}

A false-unavailable action-mask attack changes one or more availability-mask entries from available to unavailable:
\begin{equation}
    \tilde{m}_{t,j}
    =
    \begin{cases}
        0, & j \in \mathcal{S}_{\eta}^{m},\\
        m_{t,j}, & j \notin \mathcal{S}_{\eta}^{m}.
    \end{cases}
    \label{eq:action_mask_attack}
\end{equation}
This attack is a direct controller-observation attack because the action mask is consumed by the restoration policy. It is especially important when target $j$ corresponds to a load-restoration or switching action that is critical for reconnecting high-value load.

\subsubsection{Topology/Status Attack}

Let $A_t$ denote a topology or adjacency-status representation used in the graph observation. A false-open topology attack sets selected observed edges to open:
\begin{equation}
    \tilde{A}_{t,ij} = 0,
    \qquad
    (i,j) \in \mathcal{S}_{\eta}^{A},
    \label{eq:false_open_attack}
\end{equation}
whereas a false-closed hidden-outage attack sets selected observed edges to closed:
\begin{equation}
    \tilde{A}_{t,ij} = 1,
    \qquad
    (i,j) \in \mathcal{S}_{\eta}^{A}.
    \label{eq:false_closed_attack}
\end{equation}
The physical edge status in OpenDSS is not changed by \eqref{eq:false_open_attack} or \eqref{eq:false_closed_attack}; only the graph observation received by the controller is changed.

\subsubsection{Voltage Measurement Attack}

Let $v_{i,\phi,t}$ be the per-unit voltage magnitude for bus $i$ and phase $\phi$. A false-low voltage attack replaces selected active-phase voltage entries with a reported value $v^{\mathrm{low}}$:
\begin{equation}
    \tilde{v}_{i,\phi,t}
    =
    v^{\mathrm{low}},
    \qquad
    (i,\phi) \in \mathcal{S}_{\eta}^{V}.
    \label{eq:voltage_attack}
\end{equation}
The experiments use this structure to test whether the policy changes restoration decisions when the cyber observation reports artificially low voltage.

\subsubsection{Branch-Flow Hiding Attack}

Let $b_{e,t}$ be the branch-flow or branch-loading feature for edge $e$. A branch-flow hiding attack suppresses selected branch-flow observations:
\begin{equation}
    \tilde{b}_{e,t}
    =
    \alpha_b b_{e,t},
    \qquad
    e \in \mathcal{S}_{\eta}^{B},
    \label{eq:branchflow_attack}
\end{equation}
where $\alpha_b = 0$ gives full hiding. IEEE 123 results show that coordinated branch-flow hiding can be damaging even when individual branch-flow targets are weak.

\subsubsection{Global Scalar Attack}

For a global scalar feature $g_{r,t}$, an additive or multiplicative scalar attack can be written as
\begin{equation}
    \tilde{g}_{r,t}
    =
    \alpha_g g_{r,t} + \beta_g,
    \qquad
    r \in \mathcal{S}_{\eta}^{g}.
    \label{eq:global_scalar_attack}
\end{equation}
In the current evidence, these attacks are included only for scalar observations actually available to the controller.

\subsection{Clean, Attacked, and Mitigated Trajectories}
\label{subsec:trajectories}

For an outage scenario $\omega$, define the clean trajectory
\begin{equation}
    \tau^{0}(\omega)
    =
    \left(
        x^{0}_0,a^{0}_0,x^{0}_1,\ldots,x^{0}_T
    \right),
    \qquad
    a^{0}_t = \pi_{\theta}(h(x^{0}_t)).
    \label{eq:clean_trajectory}
\end{equation}
The attacked trajectory is
\begin{equation}
    \tau^{\eta}(\omega)
    =
    \left(
        x^{\eta}_0,a^{\eta}_0,x^{\eta}_1,\ldots,x^{\eta}_T
    \right),
    \qquad
    a^{\eta}_t = \pi_{\theta}(\mathcal{A}_{\eta}(h(x^{\eta}_t))).
    \label{eq:attacked_trajectory}
\end{equation}
The mitigated trajectory is
\begin{equation}
    \tau^{M}(\omega)
    =
    \left(
        x^{M}_0,a^{M}_0,x^{M}_1,\ldots,x^{M}_T
    \right),
    \qquad
    a^{M}_t = \mu_{\psi}(h(x^{M}_t),\tilde{o}^{M}_t,d_t),
    \label{eq:mitigated_trajectory}
\end{equation}
where $\mu_{\psi}$ is the defense action-selection rule, $\tilde{o}^{M}_t$ is the potentially attacked or sanitized observation available to the mitigation layer, and $d_t \in \{0,1\}$ is the detector decision. In the experiments, the clean, attacked, and mitigated trajectories for the same case are evaluated under the same outage-edge scenario and loadshape schedule.

\subsection{Physical Degradation and Recovery Metrics}
\label{subsec:metrics}

Let
\begin{equation}
    E^0 = E(\tau^0), \qquad
    E^{\eta} = E(\tau^{\eta}), \qquad
    E^{M} = E(\tau^{M})
\end{equation}
be the clean, attacked, and mitigated restored energies for the same scenario. The restored-energy loss caused by the attack is
\begin{equation}
    L^{\eta} = E^0 - E^{\eta}.
    \label{eq:loss_raw}
\end{equation}
Because an attack can occasionally improve the restoration outcome by chance, the paper reports positive physical degradation as
\begin{equation}
    D^{\eta}
    =
    \max\left(0, E^0 - E^{\eta}\right).
    \label{eq:physical_degradation}
\end{equation}
The corresponding accidental improvement is
\begin{equation}
    I^{\eta}
    =
    \max\left(0, E^{\eta} - E^0\right).
    \label{eq:accidental_improvement}
\end{equation}

Residual degradation after mitigation is
\begin{equation}
    R^{M}
    =
    \max\left(0, E^0 - E^{M}\right).
    \label{eq:residual_degradation}
\end{equation}
The degradation recovered by mitigation is
\begin{equation}
    G^{M}
    =
    D^{\eta} - R^{M}.
    \label{eq:degradation_recovered}
\end{equation}
For $D^{\eta} > \epsilon$, the clipped recovery ratio is
\begin{equation}
    \rho^{M}
    =
    \operatorname{clip}
    \left(
        \frac{D^{\eta}-R^{M}}{D^{\eta}},
        0,
        1
    \right),
    \label{eq:recovery_ratio}
\end{equation}
and $\rho^{M}=0$ when $D^{\eta}\leq \epsilon$. The clipping prevents accidental over-recovery or numerical artifacts from being interpreted as more than full recovery.

For a set of cases $\Omega$ and attack family $\eta$, the mean degradation, total degradation, and damaging-case rate are
\begin{align}
    \bar{D}_{\eta}
    &=
    \frac{1}{|\Omega|}
    \sum_{\omega \in \Omega}
    D^{\eta}(\omega),\\
    D_{\eta}^{\mathrm{tot}}
    &=
    \sum_{\omega \in \Omega}
    D^{\eta}(\omega),\\
    r_{\eta}^{\mathrm{dmg}}
    &=
    \frac{1}{|\Omega|}
    \sum_{\omega \in \Omega}
    \mathbf{1}
    \left[
        D^{\eta}(\omega) > \epsilon
    \right].
    \label{eq:aggregate_metrics}
\end{align}

\subsection{Detection Model}
\label{subsec:detection_model}

The detector maps an observation and optional reference information to a binary decision:
\begin{equation}
    d_t = \mathcal{D}_{\gamma}(\tilde{o}_t,\mathcal{R}_t)
    \in \{0,1\},
    \label{eq:detector}
\end{equation}
where $\gamma$ denotes the detector profile and $\mathcal{R}_t$ denotes available reference or consistency information. A generic residual detector computes channel residuals
\begin{equation}
    r_{k,t}
    =
    \left|
        \tilde{y}_{k,t}
        -
        \hat{y}_{k,t}
    \right|,
    \qquad
    k \in \mathcal{K},
    \label{eq:residuals}
\end{equation}
where $\tilde{y}_{k,t}$ is an observed channel and $\hat{y}_{k,t}$ is a reference, predicted, or consistency-implied value. The normalized detector score can be written as
\begin{equation}
    s_t
    =
    \max_{k \in \mathcal{K}}
    \frac{r_{k,t}}{\tau_k},
    \label{eq:detector_score}
\end{equation}
with thresholds $\tau_k>0$. The residual detector flags when
\begin{equation}
    d_t^{\mathrm{res}}
    =
    \mathbf{1}
    \left[
        s_t > 1
    \right].
    \label{eq:residual_detector}
\end{equation}

Status-integrity checks are used for binary action/status and topology/status observations. Let $\mathcal{B}$ be a set of binary channels and $\mathcal{C}_{\mathrm{stat}}$ be the set of expected logical consistency constraints among status, topology, and action-mask entries. A generic status-integrity detector is
\begin{equation}
    d_t^{\mathrm{stat}}
    =
    \mathbf{1}
    \left[
        \exists c \in \mathcal{C}_{\mathrm{stat}}
        :
        c(\tilde{o}_t) = 0
    \right].
    \label{eq:status_detector}
\end{equation}
The strict profile uses stronger status-integrity assumptions, while the relaxed profile weakens or removes selected binary-status checks to stress-test stealth behavior.

The deployed detector decision is represented as
\begin{equation}
    d_t
    =
    d_t^{\mathrm{res}}
    \lor
    d_t^{\mathrm{stat}}
    \lor
    d_t^{\mathrm{learn}},
    \label{eq:fused_detector}
\end{equation}
where $d_t^{\mathrm{learn}}$ is included only in learning-detector ablations. The learning detector does not replace the physical residual/status checks in the main evidence; it is evaluated as an additional detection signal.

For a case-level attack, detection is reported as
\begin{equation}
    d(\omega)
    =
    \mathbf{1}
    \left[
        \max_{t=0,\ldots,T-1} d_t = 1
    \right].
    \label{eq:case_detection}
\end{equation}
The detection rate for attack family $\eta$ is
\begin{equation}
    r_{\eta}^{\mathrm{det}}
    =
    \frac{1}{|\Omega|}
    \sum_{\omega \in \Omega}
    d(\omega).
    \label{eq:detection_rate}
\end{equation}
Detection on damaging cases is computed by restricting the denominator to cases with $D^{\eta}(\omega)>\epsilon$.

\subsection{Bypass and Stealth Metrics}
\label{subsec:bypass_model}

A high-impact bypass occurs when an attack is physically damaging but is not detected:
\begin{equation}
    b^{\eta}(\omega)
    =
    \mathbf{1}
    \left[
        D^{\eta}(\omega) > \epsilon_D
        \;\land\;
        d(\omega)=0
    \right],
    \label{eq:bypass_indicator}
\end{equation}
where $\epsilon_D$ is the degradation threshold used to define high-impact damage. The bypass rate is
\begin{equation}
    r_{\eta}^{\mathrm{bypass}}
    =
    \frac{1}{|\Omega|}
    \sum_{\omega \in \Omega}
    b^{\eta}(\omega).
    \label{eq:bypass_rate}
\end{equation}
For stealth-stress experiments, the same structure is used with attack-specific stealth intent. A stealth success is a case in which the attack remains below the chosen detection criterion while retaining the intended physical effect.

\subsection{Mitigation as Candidate Action Screening}
\label{subsec:mitigation_model}

After detection, mitigation is formulated as candidate action screening. Let
\begin{equation}
    \mathcal{C}_t
    =
    \left\{
        a_t^{(1)},a_t^{(2)},\ldots,a_t^{(q)}
    \right\}
    \label{eq:candidate_set}
\end{equation}
be a set of candidate restoration actions. Candidate sources may include:
\begin{align}
    a_t^{\eta} &= \pi_{\theta}(\tilde{o}_t),
    &&\text{attacked-observation policy action},\\
    a_t^{\mathrm{mit}} &= \pi_{\theta}(o_t^{\mathrm{mit}}),
    &&\text{detector-mitigated observation action},\\
    a_t^{\mathrm{ref}} &= \pi_{\theta}(o_t^{\mathrm{ref}}),
    &&\text{trusted-reference observation action}.
    \label{eq:candidate_sources}
\end{align}
The trusted-reference candidate is a robust benchmark candidate in the current evidence. It should not be described as a free replacement for detection, because it assumes access to a reference or screened observation source.

Given a candidate action $a \in \mathcal{C}_t$, define a screening score
\begin{equation}
    \widehat{Q}_t(a)
    =
    \widehat{E}_t(a)
    -
    \lambda_V \widehat{\nu}_t(a)
    -
    \lambda_S \widehat{s}_t(a),
    \label{eq:screening_score}
\end{equation}
where $\widehat{E}_t(a)$ is the estimated restored-energy contribution or rollout value, $\widehat{\nu}_t(a)$ is an estimated voltage or constraint-violation penalty, $\widehat{s}_t(a)$ is an optional switching or operational penalty, and $\lambda_V,\lambda_S \geq 0$ are weights. In the experiments, candidate quality is evaluated using the same physical scenario for clean, attacked, and mitigated comparisons; in an online deployment, \eqref{eq:screening_score} would be computed from a trusted state estimator, digital twin, or reference-screened model.

The selected mitigated action is
\begin{equation}
    a_t^{M}
    =
    \arg\max_{a \in \mathcal{C}_t}
    \widehat{Q}_t(a).
    \label{eq:mitigated_action}
\end{equation}

\subsubsection{Gated Candidate Screening}

In gated screening, mitigation is applied only if the detector gate is active:
\begin{equation}
    a_t^{M}
    =
    \begin{cases}
        \displaystyle
        \arg\max_{a \in \mathcal{C}_t^{\mathrm{gate}}}
        \widehat{Q}_t(a),
        & d_t=1,\\[2mm]
        a_t^{\eta},
        & d_t=0.
    \end{cases}
    \label{eq:gated_screening}
\end{equation}
This mode can fail when the detector misses a damaging attack or when the gated candidate set does not include a sufficiently good reference action.

\subsubsection{Always-Reference Screening}

In always-reference screening, the candidate set includes a trusted-reference action regardless of whether the detector gate is active:
\begin{equation}
    \mathcal{C}_t^{\mathrm{ref}}
    =
    \left\{
        a_t^{\eta},
        a_t^{\mathrm{mit}},
        a_t^{\mathrm{ref}}
    \right\},
    \qquad
    a_t^{M}
    =
    \arg\max_{a \in \mathcal{C}_t^{\mathrm{ref}}}
    \widehat{Q}_t(a).
    \label{eq:always_reference_screening}
\end{equation}
The IEEE 34 evidence uses this as the strongest robust-mitigation benchmark for high-impact action/status attacks.

\subsubsection{Horizon-Screened Mitigation}

For IEEE 123 advanced mitigation, candidate screening can be extended over a horizon. Let $\hat{x}_{t+h}^{a}$ be the predicted rollout state when candidate action $a$ is selected and subsequent actions follow a specified rollout rule. The horizon-screening value is
\begin{equation}
    \widehat{Q}_{t:t+H}(a)
    =
    \sum_{h=0}^{H-1}
    \left(
        \Delta t
        \sum_{i \in \mathcal{L}}
        \hat{p}^{\mathrm{srv}}_{i,t+h}(\hat{x}_{t+h}^{a})
        -
        \lambda_V
        \widehat{\nu}_{t+h}(\hat{x}_{t+h}^{a})
    \right),
    \label{eq:horizon_screening_score}
\end{equation}
and the horizon-screened action is
\begin{equation}
    a_t^{M,H}
    =
    \arg\max_{a \in \mathcal{C}_t}
    \widehat{Q}_{t:t+H}(a).
    \label{eq:horizon_screening_action}
\end{equation}
The variants `horizon\_screen\_raw\_mitigated' and `horizon\_screen\_raw\_mitigated\_reference' differ in whether the candidate pool includes only raw and mitigated actions or also trusted-reference actions.

\subsection{Observation Criticality and Coordinated Attacks}
\label{subsec:criticality_model}

For an observation target $j$ and attack family $\eta$, the single-target criticality score is the mean physical degradation over deterministic cases:
\begin{equation}
    C_{\eta}(j)
    =
    \frac{1}{|\Omega_j|}
    \sum_{\omega \in \Omega_j}
    D^{\eta}_{j}(\omega),
    \label{eq:criticality_score}
\end{equation}
where $\Omega_j$ is the set of cases in which the attack is applicable to target $j$. Targets are ranked by $C_{\eta}(j)$, maximum degradation, damaging-case count, and mitigation residuals.

For a coordinated attack on target set $\mathcal{J}$,
\begin{equation}
    \tilde{o}_t
    =
    \mathcal{A}_{\eta,\mathcal{J}}(o_t),
    \qquad
    \mathcal{J} = \{j_1,j_2,\ldots,j_k\}.
    \label{eq:coordinated_attack}
\end{equation}
The coordinated degradation is
\begin{equation}
    D^{\eta}_{\mathcal{J}}(\omega)
    =
    \max
    \left(
        0,
        E^0(\omega)
        -
        E^{\eta,\mathcal{J}}(\omega)
    \right).
    \label{eq:coordinated_degradation}
\end{equation}

For the IEEE 123 branch-flow subset ablation, let $\mathcal{J}_k$ be the top-$k$ selected branch-flow target subset and let $D_{\mathrm{base}}^{\mathrm{bf}}(\omega)$ be the baseline branch-flow hiding degradation. The total baseline reproduction ratio is
\begin{equation}
    \rho_k^{\mathrm{bf}}
    =
    \frac{
        \sum_{\omega \in \Omega_{\mathrm{bf}}}
        D^{\mathrm{bf}}_{\mathcal{J}_k}(\omega)
    }{
        \sum_{\omega \in \Omega_{\mathrm{bf}}}
        D_{\mathrm{base}}^{\mathrm{bf}}(\omega)
    },
    \label{eq:branchflow_reproduction}
\end{equation}
when the denominator is positive. This metric quantifies whether baseline branch-flow vulnerability is driven by one observation or by a coordinated subset of observations.

\subsection{Scope of the Mathematical Claims}
\label{subsec:model_scope}

The model supports claims for direct attacks on observations available to the current controller. It does not, by itself, support direct load-demand FDIA, DER-availability FDIA, or explicit fault-status FDIA unless those variables are added to the controller observation $o_t$ and evaluated with the same clean/attacked/mitigated physical workflow. This distinction is important for preserving the audit boundary of the IEEE 34-bus and IEEE 123-bus evidence packages.


\section{Mathematical Model and Problem Formulation}
\label{sec:mathematical_model}

This section formalizes the distribution restoration problem, the GRL controller observation model, the observation-level false-data injection attack (FDIA) model, and the detection/mitigation metrics used in the IEEE 34-bus and IEEE 123-bus studies. The formulation is intentionally aligned with the experimental workflow: the attacker corrupts only the cyber observation received by the controller, whereas the clean, attacked, and mitigated actions are evaluated on the same physical OpenDSS feeder state.

\subsection{Distribution Feeder and Physical State}
\label{subsec:feeder_model}

Let the distribution feeder be represented by a graph
\begin{equation}
    \mathcal{G} = (\mathcal{N},\mathcal{E}),
\end{equation}
where $\mathcal{N}$ is the set of buses and $\mathcal{E}$ is the set of distribution lines, switches, transformers, and other branch elements represented in the controller graph. For an unbalanced distribution feeder, each bus $i \in \mathcal{N}$ may have an active phase set $\Phi_i \subseteq \{a,b,c\}$, and each edge $e=(i,j) \in \mathcal{E}$ may have a phase set $\Phi_e$.

The physical feeder state at restoration step $t$ is denoted by
\begin{equation}
    x_t =
    \left(
        V_t,\,
        I_t,\,
        P_t,\,
        Q_t,\,
        z_t,\,
        \ell_t,\,
        \omega
    \right),
\end{equation}
where $V_t$ is the vector of bus voltages, $I_t$ is the vector of branch-current or branch-flow quantities, $P_t$ and $Q_t$ are real and reactive power quantities, $z_t$ contains switch and topology status variables, $\ell_t$ contains load service or controllable-load status variables, and $\omega$ denotes the outage scenario, including outaged edges, loadshape information, and other exogenous conditions. In the experiments, $x_t$ is not directly optimized in closed form; it is advanced and evaluated through OpenDSS.

The restoration action at step $t$ is a binary vector
\begin{equation}
    a_t \in \{0,1\}^{n_a},
\end{equation}
where entries correspond to controllable switch and load-restoration decisions exposed by the restoration environment. Let $\mathcal{A}(x_t)$ denote the physically feasible action set under the true feeder state. Feasibility includes operational constraints such as topology consistency, unavailable components, voltage and thermal feasibility, and any environment-specific action constraints.

The physical transition is written abstractly as
\begin{equation}
    x_{t+1} = f_{\mathrm{DSS}}(x_t,a_t,\omega_t),
    \label{eq:physical_transition}
\end{equation}
where $f_{\mathrm{DSS}}(\cdot)$ denotes the OpenDSS-backed transition and power-flow evaluation. Equation \eqref{eq:physical_transition} emphasizes that the physical feeder state evolves according to the true outage scenario and selected action, not according to the attacker's corrupted observation.

\subsection{Restoration Objective and Energy Metric}
\label{subsec:restoration_objective}

Let $p^{\mathrm{srv}}_{i,t}(x_t)$ denote the real power served at load bus or load component $i$ during restoration step $t$. For a restoration horizon of $T$ steps and step duration $\Delta t$, the restored energy of a trajectory $\tau = (x_0,a_0,x_1,\ldots,x_T)$ is
\begin{equation}
    E(\tau)
    =
    \Delta t
    \sum_{t=0}^{T-1}
    \sum_{i \in \mathcal{L}}
    p^{\mathrm{srv}}_{i,t}(x_t),
    \label{eq:restored_energy}
\end{equation}
where $\mathcal{L}$ is the set of load components included in the restoration metric. The experiments report $E(\tau)$ in kWh.

A restoration policy seeks to maximize expected restored energy, possibly with penalties for voltage violations, thermal violations, invalid actions, or switching costs. A generic policy objective can be written as
\begin{equation}
    \max_{\pi}
    J(\pi)
    =
    \mathbb{E}_{\omega}
    \left[
        \sum_{t=0}^{T-1}
        r(x_t,a_t,\omega_t)
    \right],
    \qquad
    a_t \sim \pi(\cdot \mid o_t),
    \label{eq:policy_objective}
\end{equation}
where $r(\cdot)$ is the controller training reward and $o_t$ is the controller observation. The paper's evaluation metric is not the training reward itself; it is restored-energy degradation computed from \eqref{eq:restored_energy}.

\subsection{Controller Observation and GRL Policy}
\label{subsec:observation_policy}

The restoration controller does not observe the full physical state $x_t$ directly. Instead, it receives a structured observation
\begin{equation}
    o_t = h(x_t)
    =
    \left(
        X_t,\,
        B_t,\,
        g_t,\,
        m_t
    \right),
    \label{eq:observation}
\end{equation}
where $X_t$ contains node features, $B_t$ contains edge or branch features, $g_t$ contains global restoration scalars, and $m_t \in \{0,1\}^{n_a}$ is an action/status availability mask. In the studied controller, the direct observation families include bus-voltage measurements, branch-flow measurements, topology/status features, action/status masks, and global scalar features.

The action mask defines a controller-level availability set
\begin{equation}
    \mathcal{A}_m(o_t)
    =
    \left\{
        a \in \{0,1\}^{n_a}
        \;:\;
        a_j \leq m_{t,j},
        \; j=1,\ldots,n_a
    \right\},
    \label{eq:mask_feasible_set}
\end{equation}
for action components whose convention is that $m_{t,j}=0$ marks action $j$ unavailable. The exact environment action semantics may include additional constraints, but \eqref{eq:mask_feasible_set} captures why action/status FDIA can be highly consequential: a false-unavailable mask entry can remove a useful restoration action before the learned policy selects an action.

A trained GRL policy is denoted by
\begin{equation}
    a_t^{0} = \pi_{\theta}(o_t),
    \label{eq:clean_action}
\end{equation}
where $a_t^{0}$ is the clean action selected from the uncorrupted observation. In deterministic evaluation, $\pi_{\theta}$ denotes the policy's deployed action-selection rule. For companion controller studies, the same notation also applies to non-GRL policies such as MLP controllers, with a different parameterization of $\pi_{\theta}$.

\subsection{Observation-Level FDIA Model}
\label{subsec:fdia_model}

An observation-level FDIA is modeled as an attack operator
\begin{equation}
    \tilde{o}_t = \mathcal{A}_{\eta}(o_t),
    \label{eq:attack_operator}
\end{equation}
where $\eta$ specifies the attack family, target observation indices, and attack magnitude or replacement rule. The attacked action is
\begin{equation}
    a_t^{\eta} = \pi_{\theta}(\tilde{o}_t).
    \label{eq:attacked_action}
\end{equation}
The physical feeder state remains $x_t$. Thus, $\tilde{o}_t$ affects the controller's selected action, but the action is evaluated on the true feeder through \eqref{eq:physical_transition}.

Let $\mathcal{S}_{\eta}$ be the set of observation channels modified by the attack. The attacks studied in the evidence packages are structured as follows.

\subsubsection{Action/Status Mask Attack}

A false-unavailable action-mask attack changes one or more availability-mask entries from available to unavailable:
\begin{equation}
    \tilde{m}_{t,j}
    =
    \begin{cases}
        0, & j \in \mathcal{S}_{\eta}^{m},\\
        m_{t,j}, & j \notin \mathcal{S}_{\eta}^{m}.
    \end{cases}
    \label{eq:action_mask_attack}
\end{equation}
This attack is a direct controller-observation attack because the action mask is consumed by the restoration policy. It is especially important when target $j$ corresponds to a load-restoration or switching action that is critical for reconnecting high-value load.

\subsubsection{Topology/Status Attack}

Let $A_t$ denote a topology or adjacency-status representation used in the graph observation. A false-open topology attack sets selected observed edges to open:
\begin{equation}
    \tilde{A}_{t,ij} = 0,
    \qquad
    (i,j) \in \mathcal{S}_{\eta}^{A},
    \label{eq:false_open_attack}
\end{equation}
whereas a false-closed hidden-outage attack sets selected observed edges to closed:
\begin{equation}
    \tilde{A}_{t,ij} = 1,
    \qquad
    (i,j) \in \mathcal{S}_{\eta}^{A}.
    \label{eq:false_closed_attack}
\end{equation}
The physical edge status in OpenDSS is not changed by \eqref{eq:false_open_attack} or \eqref{eq:false_closed_attack}; only the graph observation received by the controller is changed.

\subsubsection{Voltage Measurement Attack}

Let $v_{i,\phi,t}$ be the per-unit voltage magnitude for bus $i$ and phase $\phi$. A false-low voltage attack replaces selected active-phase voltage entries with a reported value $v^{\mathrm{low}}$:
\begin{equation}
    \tilde{v}_{i,\phi,t}
    =
    v^{\mathrm{low}},
    \qquad
    (i,\phi) \in \mathcal{S}_{\eta}^{V}.
    \label{eq:voltage_attack}
\end{equation}
The experiments use this structure to test whether the policy changes restoration decisions when the cyber observation reports artificially low voltage.

\subsubsection{Branch-Flow Hiding Attack}

Let $b_{e,t}$ be the branch-flow or branch-loading feature for edge $e$. A branch-flow hiding attack suppresses selected branch-flow observations:
\begin{equation}
    \tilde{b}_{e,t}
    =
    \alpha_b b_{e,t},
    \qquad
    e \in \mathcal{S}_{\eta}^{B},
    \label{eq:branchflow_attack}
\end{equation}
where $\alpha_b = 0$ gives full hiding. IEEE 123 results show that coordinated branch-flow hiding can be damaging even when individual branch-flow targets are weak.

\subsubsection{Global Scalar Attack}

For a global scalar feature $g_{r,t}$, an additive or multiplicative scalar attack can be written as
\begin{equation}
    \tilde{g}_{r,t}
    =
    \alpha_g g_{r,t} + \beta_g,
    \qquad
    r \in \mathcal{S}_{\eta}^{g}.
    \label{eq:global_scalar_attack}
\end{equation}
In the current evidence, these attacks are included only for scalar observations actually available to the controller.

\subsection{Clean, Attacked, and Mitigated Trajectories}
\label{subsec:trajectories}

For an outage scenario $\omega$, define the clean trajectory
\begin{equation}
    \tau^{0}(\omega)
    =
    \left(
        x^{0}_0,a^{0}_0,x^{0}_1,\ldots,x^{0}_T
    \right),
    \qquad
    a^{0}_t = \pi_{\theta}(h(x^{0}_t)).
    \label{eq:clean_trajectory}
\end{equation}
The attacked trajectory is
\begin{equation}
    \tau^{\eta}(\omega)
    =
    \left(
        x^{\eta}_0,a^{\eta}_0,x^{\eta}_1,\ldots,x^{\eta}_T
    \right),
    \qquad
    a^{\eta}_t = \pi_{\theta}(\mathcal{A}_{\eta}(h(x^{\eta}_t))).
    \label{eq:attacked_trajectory}
\end{equation}
The mitigated trajectory is
\begin{equation}
    \tau^{M}(\omega)
    =
    \left(
        x^{M}_0,a^{M}_0,x^{M}_1,\ldots,x^{M}_T
    \right),
    \qquad
    a^{M}_t = \mu_{\psi}(h(x^{M}_t),\tilde{o}^{M}_t,d_t),
    \label{eq:mitigated_trajectory}
\end{equation}
where $\mu_{\psi}$ is the defense action-selection rule, $\tilde{o}^{M}_t$ is the potentially attacked or sanitized observation available to the mitigation layer, and $d_t \in \{0,1\}$ is the detector decision. In the experiments, the clean, attacked, and mitigated trajectories for the same case are evaluated under the same outage-edge scenario and loadshape schedule.

\subsection{Physical Degradation and Recovery Metrics}
\label{subsec:metrics}

Let
\begin{equation}
    E^0 = E(\tau^0), \qquad
    E^{\eta} = E(\tau^{\eta}), \qquad
    E^{M} = E(\tau^{M})
\end{equation}
be the clean, attacked, and mitigated restored energies for the same scenario. The restored-energy loss caused by the attack is
\begin{equation}
    L^{\eta} = E^0 - E^{\eta}.
    \label{eq:loss_raw}
\end{equation}
Because an attack can occasionally improve the restoration outcome by chance, the paper reports positive physical degradation as
\begin{equation}
    D^{\eta}
    =
    \max\left(0, E^0 - E^{\eta}\right).
    \label{eq:physical_degradation}
\end{equation}
The corresponding accidental improvement is
\begin{equation}
    I^{\eta}
    =
    \max\left(0, E^{\eta} - E^0\right).
    \label{eq:accidental_improvement}
\end{equation}

Residual degradation after mitigation is
\begin{equation}
    R^{M}
    =
    \max\left(0, E^0 - E^{M}\right).
    \label{eq:residual_degradation}
\end{equation}
The degradation recovered by mitigation is
\begin{equation}
    G^{M}
    =
    D^{\eta} - R^{M}.
    \label{eq:degradation_recovered}
\end{equation}
For $D^{\eta} > \epsilon$, the clipped recovery ratio is
\begin{equation}
    \rho^{M}
    =
    \operatorname{clip}
    \left(
        \frac{D^{\eta}-R^{M}}{D^{\eta}},
        0,
        1
    \right),
    \label{eq:recovery_ratio}
\end{equation}
and $\rho^{M}=0$ when $D^{\eta}\leq \epsilon$. The clipping prevents accidental over-recovery or numerical artifacts from being interpreted as more than full recovery.

For a set of cases $\Omega$ and attack family $\eta$, the mean degradation, total degradation, and damaging-case rate are
\begin{align}
    \bar{D}_{\eta}
    &=
    \frac{1}{|\Omega|}
    \sum_{\omega \in \Omega}
    D^{\eta}(\omega),\\
    D_{\eta}^{\mathrm{tot}}
    &=
    \sum_{\omega \in \Omega}
    D^{\eta}(\omega),\\
    r_{\eta}^{\mathrm{dmg}}
    &=
    \frac{1}{|\Omega|}
    \sum_{\omega \in \Omega}
    \mathbf{1}
    \left[
        D^{\eta}(\omega) > \epsilon
    \right].
    \label{eq:aggregate_metrics}
\end{align}

\subsection{Detection Model}
\label{subsec:detection_model}

The detector maps an observation and optional reference information to a binary decision:
\begin{equation}
    d_t = \mathcal{D}_{\gamma}(\tilde{o}_t,\mathcal{R}_t)
    \in \{0,1\},
    \label{eq:detector}
\end{equation}
where $\gamma$ denotes the detector profile and $\mathcal{R}_t$ denotes available reference or consistency information. A generic residual detector computes channel residuals
\begin{equation}
    r_{k,t}
    =
    \left|
        \tilde{y}_{k,t}
        -
        \hat{y}_{k,t}
    \right|,
    \qquad
    k \in \mathcal{K},
    \label{eq:residuals}
\end{equation}
where $\tilde{y}_{k,t}$ is an observed channel and $\hat{y}_{k,t}$ is a reference, predicted, or consistency-implied value. The normalized detector score can be written as
\begin{equation}
    s_t
    =
    \max_{k \in \mathcal{K}}
    \frac{r_{k,t}}{\tau_k},
    \label{eq:detector_score}
\end{equation}
with thresholds $\tau_k>0$. The residual detector flags when
\begin{equation}
    d_t^{\mathrm{res}}
    =
    \mathbf{1}
    \left[
        s_t > 1
    \right].
    \label{eq:residual_detector}
\end{equation}

Status-integrity checks are used for binary action/status and topology/status observations. Let $\mathcal{B}$ be a set of binary channels and $\mathcal{C}_{\mathrm{stat}}$ be the set of expected logical consistency constraints among status, topology, and action-mask entries. A generic status-integrity detector is
\begin{equation}
    d_t^{\mathrm{stat}}
    =
    \mathbf{1}
    \left[
        \exists c \in \mathcal{C}_{\mathrm{stat}}
        :
        c(\tilde{o}_t) = 0
    \right].
    \label{eq:status_detector}
\end{equation}
The strict profile uses stronger status-integrity assumptions, while the relaxed profile weakens or removes selected binary-status checks to stress-test stealth behavior.

The deployed detector decision is represented as
\begin{equation}
    d_t
    =
    d_t^{\mathrm{res}}
    \lor
    d_t^{\mathrm{stat}}
    \lor
    d_t^{\mathrm{learn}},
    \label{eq:fused_detector}
\end{equation}
where $d_t^{\mathrm{learn}}$ is included only in learning-detector ablations. The learning detector does not replace the physical residual/status checks in the main evidence; it is evaluated as an additional detection signal.

For a case-level attack, detection is reported as
\begin{equation}
    d(\omega)
    =
    \mathbf{1}
    \left[
        \max_{t=0,\ldots,T-1} d_t = 1
    \right].
    \label{eq:case_detection}
\end{equation}
The detection rate for attack family $\eta$ is
\begin{equation}
    r_{\eta}^{\mathrm{det}}
    =
    \frac{1}{|\Omega|}
    \sum_{\omega \in \Omega}
    d(\omega).
    \label{eq:detection_rate}
\end{equation}
Detection on damaging cases is computed by restricting the denominator to cases with $D^{\eta}(\omega)>\epsilon$.

\subsection{Bypass and Stealth Metrics}
\label{subsec:bypass_model}

A high-impact bypass occurs when an attack is physically damaging but is not detected:
\begin{equation}
    b^{\eta}(\omega)
    =
    \mathbf{1}
    \left[
        D^{\eta}(\omega) > \epsilon_D
        \;\land\;
        d(\omega)=0
    \right],
    \label{eq:bypass_indicator}
\end{equation}
where $\epsilon_D$ is the degradation threshold used to define high-impact damage. The bypass rate is
\begin{equation}
    r_{\eta}^{\mathrm{bypass}}
    =
    \frac{1}{|\Omega|}
    \sum_{\omega \in \Omega}
    b^{\eta}(\omega).
    \label{eq:bypass_rate}
\end{equation}
For stealth-stress experiments, the same structure is used with attack-specific stealth intent. A stealth success is a case in which the attack remains below the chosen detection criterion while retaining the intended physical effect.

\subsection{Mitigation as Candidate Action Screening}
\label{subsec:mitigation_model}

After detection, mitigation is formulated as candidate action screening. Let
\begin{equation}
    \mathcal{C}_t
    =
    \left\{
        a_t^{(1)},a_t^{(2)},\ldots,a_t^{(q)}
    \right\}
    \label{eq:candidate_set}
\end{equation}
be a set of candidate restoration actions. Candidate sources may include:
\begin{align}
    a_t^{\eta} &= \pi_{\theta}(\tilde{o}_t),
    &&\text{attacked-observation policy action},\\
    a_t^{\mathrm{mit}} &= \pi_{\theta}(o_t^{\mathrm{mit}}),
    &&\text{detector-mitigated observation action},\\
    a_t^{\mathrm{ref}} &= \pi_{\theta}(o_t^{\mathrm{ref}}),
    &&\text{trusted-reference observation action}.
    \label{eq:candidate_sources}
\end{align}
The trusted-reference candidate is a robust benchmark candidate in the current evidence. It should not be described as a free replacement for detection, because it assumes access to a reference or screened observation source.

Given a candidate action $a \in \mathcal{C}_t$, define a screening score
\begin{equation}
    \widehat{Q}_t(a)
    =
    \widehat{E}_t(a)
    -
    \lambda_V \widehat{\nu}_t(a)
    -
    \lambda_S \widehat{s}_t(a),
    \label{eq:screening_score}
\end{equation}
where $\widehat{E}_t(a)$ is the estimated restored-energy contribution or rollout value, $\widehat{\nu}_t(a)$ is an estimated voltage or constraint-violation penalty, $\widehat{s}_t(a)$ is an optional switching or operational penalty, and $\lambda_V,\lambda_S \geq 0$ are weights. In the experiments, candidate quality is evaluated using the same physical scenario for clean, attacked, and mitigated comparisons; in an online deployment, \eqref{eq:screening_score} would be computed from a trusted state estimator, digital twin, or reference-screened model.

The selected mitigated action is
\begin{equation}
    a_t^{M}
    =
    \arg\max_{a \in \mathcal{C}_t}
    \widehat{Q}_t(a).
    \label{eq:mitigated_action}
\end{equation}

\subsubsection{Gated Candidate Screening}

In gated screening, mitigation is applied only if the detector gate is active:
\begin{equation}
    a_t^{M}
    =
    \begin{cases}
        \displaystyle
        \arg\max_{a \in \mathcal{C}_t^{\mathrm{gate}}}
        \widehat{Q}_t(a),
        & d_t=1,\\[2mm]
        a_t^{\eta},
        & d_t=0.
    \end{cases}
    \label{eq:gated_screening}
\end{equation}
This mode can fail when the detector misses a damaging attack or when the gated candidate set does not include a sufficiently good reference action.

\subsubsection{Always-Reference Screening}

In always-reference screening, the candidate set includes a trusted-reference action regardless of whether the detector gate is active:
\begin{equation}
    \mathcal{C}_t^{\mathrm{ref}}
    =
    \left\{
        a_t^{\eta},
        a_t^{\mathrm{mit}},
        a_t^{\mathrm{ref}}
    \right\},
    \qquad
    a_t^{M}
    =
    \arg\max_{a \in \mathcal{C}_t^{\mathrm{ref}}}
    \widehat{Q}_t(a).
    \label{eq:always_reference_screening}
\end{equation}
The IEEE 34 evidence uses this as the strongest robust-mitigation benchmark for high-impact action/status attacks.

\subsubsection{Horizon-Screened Mitigation}

For IEEE 123 advanced mitigation, candidate screening can be extended over a horizon. Let $\hat{x}_{t+h}^{a}$ be the predicted rollout state when candidate action $a$ is selected and subsequent actions follow a specified rollout rule. The horizon-screening value is
\begin{equation}
    \widehat{Q}_{t:t+H}(a)
    =
    \sum_{h=0}^{H-1}
    \left(
        \Delta t
        \sum_{i \in \mathcal{L}}
        \hat{p}^{\mathrm{srv}}_{i,t+h}(\hat{x}_{t+h}^{a})
        -
        \lambda_V
        \widehat{\nu}_{t+h}(\hat{x}_{t+h}^{a})
    \right),
    \label{eq:horizon_screening_score}
\end{equation}
and the horizon-screened action is
\begin{equation}
    a_t^{M,H}
    =
    \arg\max_{a \in \mathcal{C}_t}
    \widehat{Q}_{t:t+H}(a).
    \label{eq:horizon_screening_action}
\end{equation}
The variants `horizon\_screen\_raw\_mitigated' and `horizon\_screen\_raw\_mitigated\_reference' differ in whether the candidate pool includes only raw and mitigated actions or also trusted-reference actions.

\subsection{Observation Criticality and Coordinated Attacks}
\label{subsec:criticality_model}

For an observation target $j$ and attack family $\eta$, the single-target criticality score is the mean physical degradation over deterministic cases:
\begin{equation}
    C_{\eta}(j)
    =
    \frac{1}{|\Omega_j|}
    \sum_{\omega \in \Omega_j}
    D^{\eta}_{j}(\omega),
    \label{eq:criticality_score}
\end{equation}
where $\Omega_j$ is the set of cases in which the attack is applicable to target $j$. Targets are ranked by $C_{\eta}(j)$, maximum degradation, damaging-case count, and mitigation residuals.

For a coordinated attack on target set $\mathcal{J}$,
\begin{equation}
    \tilde{o}_t
    =
    \mathcal{A}_{\eta,\mathcal{J}}(o_t),
    \qquad
    \mathcal{J} = \{j_1,j_2,\ldots,j_k\}.
    \label{eq:coordinated_attack}
\end{equation}
The coordinated degradation is
\begin{equation}
    D^{\eta}_{\mathcal{J}}(\omega)
    =
    \max
    \left(
        0,
        E^0(\omega)
        -
        E^{\eta,\mathcal{J}}(\omega)
    \right).
    \label{eq:coordinated_degradation}
\end{equation}

For the IEEE 123 branch-flow subset ablation, let $\mathcal{J}_k$ be the top-$k$ selected branch-flow target subset and let $D_{\mathrm{base}}^{\mathrm{bf}}(\omega)$ be the baseline branch-flow hiding degradation. The total baseline reproduction ratio is
\begin{equation}
    \rho_k^{\mathrm{bf}}
    =
    \frac{
        \sum_{\omega \in \Omega_{\mathrm{bf}}}
        D^{\mathrm{bf}}_{\mathcal{J}_k}(\omega)
    }{
        \sum_{\omega \in \Omega_{\mathrm{bf}}}
        D_{\mathrm{base}}^{\mathrm{bf}}(\omega)
    },
    \label{eq:branchflow_reproduction}
\end{equation}
when the denominator is positive. This metric quantifies whether baseline branch-flow vulnerability is driven by one observation or by a coordinated subset of observations.

\subsection{Scope of the Mathematical Claims}
\label{subsec:model_scope}

The model supports claims for direct attacks on observations available to the current controller. It does not, by itself, support direct load-demand FDIA, DER-availability FDIA, or explicit fault-status FDIA unless those variables are added to the controller observation $o_t$ and evaluated with the same clean/attacked/mitigated physical workflow. This distinction is important for preserving the audit boundary of the IEEE 34-bus and IEEE 123-bus evidence packages.


\section{Mathematical Model and Problem Formulation}
\label{sec:mathematical_model}

This section formalizes the distribution restoration problem, the GRL controller observation model, the observation-level false-data injection attack (FDIA) model, and the detection/mitigation metrics used in the IEEE 34-bus and IEEE 123-bus studies. The formulation is intentionally aligned with the experimental workflow: the attacker corrupts only the cyber observation received by the controller, whereas the clean, attacked, and mitigated actions are evaluated on the same physical OpenDSS feeder state.

\subsection{Distribution Feeder and Physical State}
\label{subsec:feeder_model}

Let the distribution feeder be represented by a graph
\begin{equation}
    \mathcal{G} = (\mathcal{N},\mathcal{E}),
\end{equation}
where $\mathcal{N}$ is the set of buses and $\mathcal{E}$ is the set of distribution lines, switches, transformers, and other branch elements represented in the controller graph. For an unbalanced distribution feeder, each bus $i \in \mathcal{N}$ may have an active phase set $\Phi_i \subseteq \{a,b,c\}$, and each edge $e=(i,j) \in \mathcal{E}$ may have a phase set $\Phi_e$.

The physical feeder state at restoration step $t$ is denoted by
\begin{equation}
    x_t =
    \left(
        V_t,\,
        I_t,\,
        P_t,\,
        Q_t,\,
        z_t,\,
        \ell_t,\,
        \omega
    \right),
\end{equation}
where $V_t$ is the vector of bus voltages, $I_t$ is the vector of branch-current or branch-flow quantities, $P_t$ and $Q_t$ are real and reactive power quantities, $z_t$ contains switch and topology status variables, $\ell_t$ contains load service or controllable-load status variables, and $\omega$ denotes the outage scenario, including outaged edges, loadshape information, and other exogenous conditions. In the experiments, $x_t$ is not directly optimized in closed form; it is advanced and evaluated through OpenDSS.

The restoration action at step $t$ is a binary vector
\begin{equation}
    a_t \in \{0,1\}^{n_a},
\end{equation}
where entries correspond to controllable switch and load-restoration decisions exposed by the restoration environment. Let $\mathcal{A}(x_t)$ denote the physically feasible action set under the true feeder state. Feasibility includes operational constraints such as topology consistency, unavailable components, voltage and thermal feasibility, and any environment-specific action constraints.

The physical transition is written abstractly as
\begin{equation}
    x_{t+1} = f_{\mathrm{DSS}}(x_t,a_t,\omega_t),
    \label{eq:physical_transition}
\end{equation}
where $f_{\mathrm{DSS}}(\cdot)$ denotes the OpenDSS-backed transition and power-flow evaluation. Equation \eqref{eq:physical_transition} emphasizes that the physical feeder state evolves according to the true outage scenario and selected action, not according to the attacker's corrupted observation.

\subsection{Restoration Objective and Energy Metric}
\label{subsec:restoration_objective}

Let $p^{\mathrm{srv}}_{i,t}(x_t)$ denote the real power served at load bus or load component $i$ during restoration step $t$. For a restoration horizon of $T$ steps and step duration $\Delta t$, the restored energy of a trajectory $\tau = (x_0,a_0,x_1,\ldots,x_T)$ is
\begin{equation}
    E(\tau)
    =
    \Delta t
    \sum_{t=0}^{T-1}
    \sum_{i \in \mathcal{L}}
    p^{\mathrm{srv}}_{i,t}(x_t),
    \label{eq:restored_energy}
\end{equation}
where $\mathcal{L}$ is the set of load components included in the restoration metric. The experiments report $E(\tau)$ in kWh.

A restoration policy seeks to maximize expected restored energy, possibly with penalties for voltage violations, thermal violations, invalid actions, or switching costs. A generic policy objective can be written as
\begin{equation}
    \max_{\pi}
    J(\pi)
    =
    \mathbb{E}_{\omega}
    \left[
        \sum_{t=0}^{T-1}
        r(x_t,a_t,\omega_t)
    \right],
    \qquad
    a_t \sim \pi(\cdot \mid o_t),
    \label{eq:policy_objective}
\end{equation}
where $r(\cdot)$ is the controller training reward and $o_t$ is the controller observation. The paper's evaluation metric is not the training reward itself; it is restored-energy degradation computed from \eqref{eq:restored_energy}.

\subsection{Controller Observation and GRL Policy}
\label{subsec:observation_policy}

The restoration controller does not observe the full physical state $x_t$ directly. Instead, it receives a structured observation
\begin{equation}
    o_t = h(x_t)
    =
    \left(
        X_t,\,
        B_t,\,
        g_t,\,
        m_t
    \right),
    \label{eq:observation}
\end{equation}
where $X_t$ contains node features, $B_t$ contains edge or branch features, $g_t$ contains global restoration scalars, and $m_t \in \{0,1\}^{n_a}$ is an action/status availability mask. In the studied controller, the direct observation families include bus-voltage measurements, branch-flow measurements, topology/status features, action/status masks, and global scalar features.

The action mask defines a controller-level availability set
\begin{equation}
    \mathcal{A}_m(o_t)
    =
    \left\{
        a \in \{0,1\}^{n_a}
        \;:\;
        a_j \leq m_{t,j},
        \; j=1,\ldots,n_a
    \right\},
    \label{eq:mask_feasible_set}
\end{equation}
for action components whose convention is that $m_{t,j}=0$ marks action $j$ unavailable. The exact environment action semantics may include additional constraints, but \eqref{eq:mask_feasible_set} captures why action/status FDIA can be highly consequential: a false-unavailable mask entry can remove a useful restoration action before the learned policy selects an action.

A trained GRL policy is denoted by
\begin{equation}
    a_t^{0} = \pi_{\theta}(o_t),
    \label{eq:clean_action}
\end{equation}
where $a_t^{0}$ is the clean action selected from the uncorrupted observation. In deterministic evaluation, $\pi_{\theta}$ denotes the policy's deployed action-selection rule. For companion controller studies, the same notation also applies to non-GRL policies such as MLP controllers, with a different parameterization of $\pi_{\theta}$.

\subsection{Observation-Level FDIA Model}
\label{subsec:fdia_model}

An observation-level FDIA is modeled as an attack operator
\begin{equation}
    \tilde{o}_t = \mathcal{A}_{\eta}(o_t),
    \label{eq:attack_operator}
\end{equation}
where $\eta$ specifies the attack family, target observation indices, and attack magnitude or replacement rule. The attacked action is
\begin{equation}
    a_t^{\eta} = \pi_{\theta}(\tilde{o}_t).
    \label{eq:attacked_action}
\end{equation}
The physical feeder state remains $x_t$. Thus, $\tilde{o}_t$ affects the controller's selected action, but the action is evaluated on the true feeder through \eqref{eq:physical_transition}.

Let $\mathcal{S}_{\eta}$ be the set of observation channels modified by the attack. The attacks studied in the evidence packages are structured as follows.

\subsubsection{Action/Status Mask Attack}

A false-unavailable action-mask attack changes one or more availability-mask entries from available to unavailable:
\begin{equation}
    \tilde{m}_{t,j}
    =
    \begin{cases}
        0, & j \in \mathcal{S}_{\eta}^{m},\\
        m_{t,j}, & j \notin \mathcal{S}_{\eta}^{m}.
    \end{cases}
    \label{eq:action_mask_attack}
\end{equation}
This attack is a direct controller-observation attack because the action mask is consumed by the restoration policy. It is especially important when target $j$ corresponds to a load-restoration or switching action that is critical for reconnecting high-value load.

\subsubsection{Topology/Status Attack}

Let $A_t$ denote a topology or adjacency-status representation used in the graph observation. A false-open topology attack sets selected observed edges to open:
\begin{equation}
    \tilde{A}_{t,ij} = 0,
    \qquad
    (i,j) \in \mathcal{S}_{\eta}^{A},
    \label{eq:false_open_attack}
\end{equation}
whereas a false-closed hidden-outage attack sets selected observed edges to closed:
\begin{equation}
    \tilde{A}_{t,ij} = 1,
    \qquad
    (i,j) \in \mathcal{S}_{\eta}^{A}.
    \label{eq:false_closed_attack}
\end{equation}
The physical edge status in OpenDSS is not changed by \eqref{eq:false_open_attack} or \eqref{eq:false_closed_attack}; only the graph observation received by the controller is changed.

\subsubsection{Voltage Measurement Attack}

Let $v_{i,\phi,t}$ be the per-unit voltage magnitude for bus $i$ and phase $\phi$. A false-low voltage attack replaces selected active-phase voltage entries with a reported value $v^{\mathrm{low}}$:
\begin{equation}
    \tilde{v}_{i,\phi,t}
    =
    v^{\mathrm{low}},
    \qquad
    (i,\phi) \in \mathcal{S}_{\eta}^{V}.
    \label{eq:voltage_attack}
\end{equation}
The experiments use this structure to test whether the policy changes restoration decisions when the cyber observation reports artificially low voltage.

\subsubsection{Branch-Flow Hiding Attack}

Let $b_{e,t}$ be the branch-flow or branch-loading feature for edge $e$. A branch-flow hiding attack suppresses selected branch-flow observations:
\begin{equation}
    \tilde{b}_{e,t}
    =
    \alpha_b b_{e,t},
    \qquad
    e \in \mathcal{S}_{\eta}^{B},
    \label{eq:branchflow_attack}
\end{equation}
where $\alpha_b = 0$ gives full hiding. IEEE 123 results show that coordinated branch-flow hiding can be damaging even when individual branch-flow targets are weak.

\subsubsection{Global Scalar Attack}

For a global scalar feature $g_{r,t}$, an additive or multiplicative scalar attack can be written as
\begin{equation}
    \tilde{g}_{r,t}
    =
    \alpha_g g_{r,t} + \beta_g,
    \qquad
    r \in \mathcal{S}_{\eta}^{g}.
    \label{eq:global_scalar_attack}
\end{equation}
In the current evidence, these attacks are included only for scalar observations actually available to the controller.

\subsection{Clean, Attacked, and Mitigated Trajectories}
\label{subsec:trajectories}

For an outage scenario $\omega$, define the clean trajectory
\begin{equation}
    \tau^{0}(\omega)
    =
    \left(
        x^{0}_0,a^{0}_0,x^{0}_1,\ldots,x^{0}_T
    \right),
    \qquad
    a^{0}_t = \pi_{\theta}(h(x^{0}_t)).
    \label{eq:clean_trajectory}
\end{equation}
The attacked trajectory is
\begin{equation}
    \tau^{\eta}(\omega)
    =
    \left(
        x^{\eta}_0,a^{\eta}_0,x^{\eta}_1,\ldots,x^{\eta}_T
    \right),
    \qquad
    a^{\eta}_t = \pi_{\theta}(\mathcal{A}_{\eta}(h(x^{\eta}_t))).
    \label{eq:attacked_trajectory}
\end{equation}
The mitigated trajectory is
\begin{equation}
    \tau^{M}(\omega)
    =
    \left(
        x^{M}_0,a^{M}_0,x^{M}_1,\ldots,x^{M}_T
    \right),
    \qquad
    a^{M}_t = \mu_{\psi}(h(x^{M}_t),\tilde{o}^{M}_t,d_t),
    \label{eq:mitigated_trajectory}
\end{equation}
where $\mu_{\psi}$ is the defense action-selection rule, $\tilde{o}^{M}_t$ is the potentially attacked or sanitized observation available to the mitigation layer, and $d_t \in \{0,1\}$ is the detector decision. In the experiments, the clean, attacked, and mitigated trajectories for the same case are evaluated under the same outage-edge scenario and loadshape schedule.

\subsection{Physical Degradation and Recovery Metrics}
\label{subsec:metrics}

Let
\begin{equation}
    E^0 = E(\tau^0), \qquad
    E^{\eta} = E(\tau^{\eta}), \qquad
    E^{M} = E(\tau^{M})
\end{equation}
be the clean, attacked, and mitigated restored energies for the same scenario. The restored-energy loss caused by the attack is
\begin{equation}
    L^{\eta} = E^0 - E^{\eta}.
    \label{eq:loss_raw}
\end{equation}
Because an attack can occasionally improve the restoration outcome by chance, the paper reports positive physical degradation as
\begin{equation}
    D^{\eta}
    =
    \max\left(0, E^0 - E^{\eta}\right).
    \label{eq:physical_degradation}
\end{equation}
The corresponding accidental improvement is
\begin{equation}
    I^{\eta}
    =
    \max\left(0, E^{\eta} - E^0\right).
    \label{eq:accidental_improvement}
\end{equation}

Residual degradation after mitigation is
\begin{equation}
    R^{M}
    =
    \max\left(0, E^0 - E^{M}\right).
    \label{eq:residual_degradation}
\end{equation}
The degradation recovered by mitigation is
\begin{equation}
    G^{M}
    =
    D^{\eta} - R^{M}.
    \label{eq:degradation_recovered}
\end{equation}
For $D^{\eta} > \epsilon$, the clipped recovery ratio is
\begin{equation}
    \rho^{M}
    =
    \operatorname{clip}
    \left(
        \frac{D^{\eta}-R^{M}}{D^{\eta}},
        0,
        1
    \right),
    \label{eq:recovery_ratio}
\end{equation}
and $\rho^{M}=0$ when $D^{\eta}\leq \epsilon$. The clipping prevents accidental over-recovery or numerical artifacts from being interpreted as more than full recovery.

For a set of cases $\Omega$ and attack family $\eta$, the mean degradation, total degradation, and damaging-case rate are
\begin{align}
    \bar{D}_{\eta}
    &=
    \frac{1}{|\Omega|}
    \sum_{\omega \in \Omega}
    D^{\eta}(\omega),\\
    D_{\eta}^{\mathrm{tot}}
    &=
    \sum_{\omega \in \Omega}
    D^{\eta}(\omega),\\
    r_{\eta}^{\mathrm{dmg}}
    &=
    \frac{1}{|\Omega|}
    \sum_{\omega \in \Omega}
    \mathbf{1}
    \left[
        D^{\eta}(\omega) > \epsilon
    \right].
    \label{eq:aggregate_metrics}
\end{align}

\subsection{Detection Model}
\label{subsec:detection_model}

The detector maps an observation and optional reference information to a binary decision:
\begin{equation}
    d_t = \mathcal{D}_{\gamma}(\tilde{o}_t,\mathcal{R}_t)
    \in \{0,1\},
    \label{eq:detector}
\end{equation}
where $\gamma$ denotes the detector profile and $\mathcal{R}_t$ denotes available reference or consistency information. A generic residual detector computes channel residuals
\begin{equation}
    r_{k,t}
    =
    \left|
        \tilde{y}_{k,t}
        -
        \hat{y}_{k,t}
    \right|,
    \qquad
    k \in \mathcal{K},
    \label{eq:residuals}
\end{equation}
where $\tilde{y}_{k,t}$ is an observed channel and $\hat{y}_{k,t}$ is a reference, predicted, or consistency-implied value. The normalized detector score can be written as
\begin{equation}
    s_t
    =
    \max_{k \in \mathcal{K}}
    \frac{r_{k,t}}{\tau_k},
    \label{eq:detector_score}
\end{equation}
with thresholds $\tau_k>0$. The residual detector flags when
\begin{equation}
    d_t^{\mathrm{res}}
    =
    \mathbf{1}
    \left[
        s_t > 1
    \right].
    \label{eq:residual_detector}
\end{equation}

Status-integrity checks are used for binary action/status and topology/status observations. Let $\mathcal{B}$ be a set of binary channels and $\mathcal{C}_{\mathrm{stat}}$ be the set of expected logical consistency constraints among status, topology, and action-mask entries. A generic status-integrity detector is
\begin{equation}
    d_t^{\mathrm{stat}}
    =
    \mathbf{1}
    \left[
        \exists c \in \mathcal{C}_{\mathrm{stat}}
        :
        c(\tilde{o}_t) = 0
    \right].
    \label{eq:status_detector}
\end{equation}
The strict profile uses stronger status-integrity assumptions, while the relaxed profile weakens or removes selected binary-status checks to stress-test stealth behavior.

The deployed detector decision is represented as
\begin{equation}
    d_t
    =
    d_t^{\mathrm{res}}
    \lor
    d_t^{\mathrm{stat}}
    \lor
    d_t^{\mathrm{learn}},
    \label{eq:fused_detector}
\end{equation}
where $d_t^{\mathrm{learn}}$ is included only in learning-detector ablations. The learning detector does not replace the physical residual/status checks in the main evidence; it is evaluated as an additional detection signal.

For a case-level attack, detection is reported as
\begin{equation}
    d(\omega)
    =
    \mathbf{1}
    \left[
        \max_{t=0,\ldots,T-1} d_t = 1
    \right].
    \label{eq:case_detection}
\end{equation}
The detection rate for attack family $\eta$ is
\begin{equation}
    r_{\eta}^{\mathrm{det}}
    =
    \frac{1}{|\Omega|}
    \sum_{\omega \in \Omega}
    d(\omega).
    \label{eq:detection_rate}
\end{equation}
Detection on damaging cases is computed by restricting the denominator to cases with $D^{\eta}(\omega)>\epsilon$.

\subsection{Bypass and Stealth Metrics}
\label{subsec:bypass_model}

A high-impact bypass occurs when an attack is physically damaging but is not detected:
\begin{equation}
    b^{\eta}(\omega)
    =
    \mathbf{1}
    \left[
        D^{\eta}(\omega) > \epsilon_D
        \;\land\;
        d(\omega)=0
    \right],
    \label{eq:bypass_indicator}
\end{equation}
where $\epsilon_D$ is the degradation threshold used to define high-impact damage. The bypass rate is
\begin{equation}
    r_{\eta}^{\mathrm{bypass}}
    =
    \frac{1}{|\Omega|}
    \sum_{\omega \in \Omega}
    b^{\eta}(\omega).
    \label{eq:bypass_rate}
\end{equation}
For stealth-stress experiments, the same structure is used with attack-specific stealth intent. A stealth success is a case in which the attack remains below the chosen detection criterion while retaining the intended physical effect.

\subsection{Mitigation as Candidate Action Screening}
\label{subsec:mitigation_model}

After detection, mitigation is formulated as candidate action screening. Let
\begin{equation}
    \mathcal{C}_t
    =
    \left\{
        a_t^{(1)},a_t^{(2)},\ldots,a_t^{(q)}
    \right\}
    \label{eq:candidate_set}
\end{equation}
be a set of candidate restoration actions. Candidate sources may include:
\begin{align}
    a_t^{\eta} &= \pi_{\theta}(\tilde{o}_t),
    &&\text{attacked-observation policy action},\\
    a_t^{\mathrm{mit}} &= \pi_{\theta}(o_t^{\mathrm{mit}}),
    &&\text{detector-mitigated observation action},\\
    a_t^{\mathrm{ref}} &= \pi_{\theta}(o_t^{\mathrm{ref}}),
    &&\text{trusted-reference observation action}.
    \label{eq:candidate_sources}
\end{align}
The trusted-reference candidate is a robust benchmark candidate in the current evidence. It should not be described as a free replacement for detection, because it assumes access to a reference or screened observation source.

Given a candidate action $a \in \mathcal{C}_t$, define a screening score
\begin{equation}
    \widehat{Q}_t(a)
    =
    \widehat{E}_t(a)
    -
    \lambda_V \widehat{\nu}_t(a)
    -
    \lambda_S \widehat{s}_t(a),
    \label{eq:screening_score}
\end{equation}
where $\widehat{E}_t(a)$ is the estimated restored-energy contribution or rollout value, $\widehat{\nu}_t(a)$ is an estimated voltage or constraint-violation penalty, $\widehat{s}_t(a)$ is an optional switching or operational penalty, and $\lambda_V,\lambda_S \geq 0$ are weights. In the experiments, candidate quality is evaluated using the same physical scenario for clean, attacked, and mitigated comparisons; in an online deployment, \eqref{eq:screening_score} would be computed from a trusted state estimator, digital twin, or reference-screened model.

The selected mitigated action is
\begin{equation}
    a_t^{M}
    =
    \arg\max_{a \in \mathcal{C}_t}
    \widehat{Q}_t(a).
    \label{eq:mitigated_action}
\end{equation}

\subsubsection{Gated Candidate Screening}

In gated screening, mitigation is applied only if the detector gate is active:
\begin{equation}
    a_t^{M}
    =
    \begin{cases}
        \displaystyle
        \arg\max_{a \in \mathcal{C}_t^{\mathrm{gate}}}
        \widehat{Q}_t(a),
        & d_t=1,\\[2mm]
        a_t^{\eta},
        & d_t=0.
    \end{cases}
    \label{eq:gated_screening}
\end{equation}
This mode can fail when the detector misses a damaging attack or when the gated candidate set does not include a sufficiently good reference action.

\subsubsection{Always-Reference Screening}

In always-reference screening, the candidate set includes a trusted-reference action regardless of whether the detector gate is active:
\begin{equation}
    \mathcal{C}_t^{\mathrm{ref}}
    =
    \left\{
        a_t^{\eta},
        a_t^{\mathrm{mit}},
        a_t^{\mathrm{ref}}
    \right\},
    \qquad
    a_t^{M}
    =
    \arg\max_{a \in \mathcal{C}_t^{\mathrm{ref}}}
    \widehat{Q}_t(a).
    \label{eq:always_reference_screening}
\end{equation}
The IEEE 34 evidence uses this as the strongest robust-mitigation benchmark for high-impact action/status attacks.

\subsubsection{Horizon-Screened Mitigation}

For IEEE 123 advanced mitigation, candidate screening can be extended over a horizon. Let $\hat{x}_{t+h}^{a}$ be the predicted rollout state when candidate action $a$ is selected and subsequent actions follow a specified rollout rule. The horizon-screening value is
\begin{equation}
    \widehat{Q}_{t:t+H}(a)
    =
    \sum_{h=0}^{H-1}
    \left(
        \Delta t
        \sum_{i \in \mathcal{L}}
        \hat{p}^{\mathrm{srv}}_{i,t+h}(\hat{x}_{t+h}^{a})
        -
        \lambda_V
        \widehat{\nu}_{t+h}(\hat{x}_{t+h}^{a})
    \right),
    \label{eq:horizon_screening_score}
\end{equation}
and the horizon-screened action is
\begin{equation}
    a_t^{M,H}
    =
    \arg\max_{a \in \mathcal{C}_t}
    \widehat{Q}_{t:t+H}(a).
    \label{eq:horizon_screening_action}
\end{equation}
The variants `horizon\_screen\_raw\_mitigated' and `horizon\_screen\_raw\_mitigated\_reference' differ in whether the candidate pool includes only raw and mitigated actions or also trusted-reference actions.

\subsection{Observation Criticality and Coordinated Attacks}
\label{subsec:criticality_model}

For an observation target $j$ and attack family $\eta$, the single-target criticality score is the mean physical degradation over deterministic cases:
\begin{equation}
    C_{\eta}(j)
    =
    \frac{1}{|\Omega_j|}
    \sum_{\omega \in \Omega_j}
    D^{\eta}_{j}(\omega),
    \label{eq:criticality_score}
\end{equation}
where $\Omega_j$ is the set of cases in which the attack is applicable to target $j$. Targets are ranked by $C_{\eta}(j)$, maximum degradation, damaging-case count, and mitigation residuals.

For a coordinated attack on target set $\mathcal{J}$,
\begin{equation}
    \tilde{o}_t
    =
    \mathcal{A}_{\eta,\mathcal{J}}(o_t),
    \qquad
    \mathcal{J} = \{j_1,j_2,\ldots,j_k\}.
    \label{eq:coordinated_attack}
\end{equation}
The coordinated degradation is
\begin{equation}
    D^{\eta}_{\mathcal{J}}(\omega)
    =
    \max
    \left(
        0,
        E^0(\omega)
        -
        E^{\eta,\mathcal{J}}(\omega)
    \right).
    \label{eq:coordinated_degradation}
\end{equation}

For the IEEE 123 branch-flow subset ablation, let $\mathcal{J}_k$ be the top-$k$ selected branch-flow target subset and let $D_{\mathrm{base}}^{\mathrm{bf}}(\omega)$ be the baseline branch-flow hiding degradation. The total baseline reproduction ratio is
\begin{equation}
    \rho_k^{\mathrm{bf}}
    =
    \frac{
        \sum_{\omega \in \Omega_{\mathrm{bf}}}
        D^{\mathrm{bf}}_{\mathcal{J}_k}(\omega)
    }{
        \sum_{\omega \in \Omega_{\mathrm{bf}}}
        D_{\mathrm{base}}^{\mathrm{bf}}(\omega)
    },
    \label{eq:branchflow_reproduction}
\end{equation}
when the denominator is positive. This metric quantifies whether baseline branch-flow vulnerability is driven by one observation or by a coordinated subset of observations.

\subsection{Scope of the Mathematical Claims}
\label{subsec:model_scope}

The model supports claims for direct attacks on observations available to the current controller. It does not, by itself, support direct load-demand FDIA, DER-availability FDIA, or explicit fault-status FDIA unless those variables are added to the controller observation $o_t$ and evaluated with the same clean/attacked/mitigated physical workflow. This distinction is important for preserving the audit boundary of the IEEE 34-bus and IEEE 123-bus evidence packages.
